\chapter{System Architecture and Design}
\label{ch:architecture}

\section{Layered architecture}

The four layers (\cref{fig:thesis-architecture}) are a React route UI; FastAPI public-projection/protected APIs and security; ingestion, retrieval, selective QA, structured recommendation, conversational idea generation, discovery, review and evaluation logic; and SQLite plus immutable local generations. Production routes use the approved public projection. Only protected evaluation or an explicitly labelled loopback demo may select technically eligible records.

\begin{figure}[htbp]
\centering
\resizebox{0.96\textwidth}{!}{\begin{tikzpicture}[node distance=5mm]
  \node[layerbox] (ui) {\textbf{Presentation layer} \hfill React 19 + Vite + TypeScript\\Dashboard; Paper Browser/Detail; Search; Ask; conversational Idea Generator; Explorer; Evaluation; protected Admin Control; provenance, model, freshness, privacy and review notices};
  \node[layerbox, below=of ui, fill=PaleTeal, draw=Teal!80!black] (api) {\textbf{API and security layer} \hfill FastAPI\\Production public projection and transient generation; session/role-protected mutations plus service actors; CSRF; feature switches; input/rate limits; exact CORS/hosts; redacted logs; readiness};
  \node[layerbox, below=of api, fill=PaleGold, draw=Gold!80!black] (logic) {\textbf{Research-intelligence and evaluation layer}\\Discovery, safe PDF acquisition, native/OCR-aware parsing and chunking; keyword, feature-hashing, pinned learned-dense and hybrid retrieval; QA; evaluated structured Finder; source-aware Idea Generator; controlled topics; artefacts; offline evaluation};
  \node[layerbox, below=of logic, fill=gray!8, draw=Slate] (data) {\textbf{Persistence and evidence layer} \hfill SQLite + immutable generation directories\\Source/provenance entities; permitted PDFs and derived text; current pointers plus hash-validated manifests; silver labels; raw rankings, claims and reviews};
  \draw[arrow] (ui) -- node[right,font=\scriptsize]{JSON/HTTP} (api);
  \draw[arrow] (api) -- (logic);
  \draw[arrow] (logic) -- node[right,font=\scriptsize]{SQLModel + atomic files} (data);
  \draw[dashedarrow] (data.west) .. controls +(-16mm,0) and +(-16mm,0) .. node[left,font=\scriptsize,align=center]{paper/chunk/page/section,\\provider, timestamp, warnings, review} (ui.west);
\end{tikzpicture}
}
\caption[Layered system architecture and provenance return path.]{Layered system architecture. The dashed path represents provenance and review data returned from storage to public evidence displays. Text labels provide the non-color encoding.}
\label{fig:thesis-architecture}
\end{figure}

\section{Design principles}

Seven principles govern the design. First, derived index state is authoritative only when its manifest and active generation pointer match the current corpus. Second, technical eligibility never implies publication or redistribution approval. Third, source facts, paper-informed ideas, and general suggestions use distinct fields and labels. Fourth, AI review has a distinct \texttt{ai\_reviewed} state and cannot create human approval. Fifth, public requests are privacy-minimised and administrative mutations are authenticated. Sixth, model and retrieval identity are recorded at generation time. Seventh, missing dependencies or evidence produce explicit unavailable, stale, partial, \texttt{support\_unverified}, or needs-reprocess states instead of silent success.

\section{Persistence and traceability model}

\Cref{fig:entity-model} summarises the entities. Papers/chunks store extraction, chunk and approved-public generations; outputs retain structured source identifiers. Controlled topics link papers and canonical identities. Each transition appends an actor-typed, request-attributed, hash-linked \texttt{ReviewEvent}, preserving history without treating AI review as human approval.

\begin{figure}[htbp]
\centering
\resizebox{0.97\textwidth}{!}{\begin{tikzpicture}[node distance=9mm and 11mm]
  \node[entity] (paper) {\textbf{Paper}\\\underline{paper\_id}\\title, authors JSON\\year, venue, URLs\\extraction/review status};
  \node[entity, right=of paper] (chunk) {\textbf{Chunk}\\\underline{chunk\_id}\\paper\_id (FK)\\pages, section, text\\hash, embedding status};
  \node[entity, right=of chunk] (answer) {\textbf{RAGAnswer}\\\underline{answer\_id}\\question, answer\\chunk/paper/citation JSON\\provider, grounding, review};

  \node[entity, below=of paper] (artifact) {\textbf{PaperArtifact}\\\underline{artifact\_id}\\paper\_id (FK)\\generated/citation JSON\\grounding/review status};
  \node[entity, below=of chunk] (recommendation) {\textbf{ThesisRecommendation}\\\underline{recommendation\_id}\\student request JSON\\recommendations/citations\\grounding/review status};
  \node[entity, below=of answer] (review) {\textbf{ReviewEvent}\\\underline{review\_event\_id}\\item\_type + item\_id\\actor type/role + request ID\\status, diff, prior/event hashes};

  \node[entity, below=12mm of artifact] (topic) {\textbf{Topic}\\\underline{topic\_id}\\normalized name\\source, review status};
  \node[entity, right=of topic] (papertopic) {\textbf{PaperTopic}\\\underline{link\_id}\\paper\_id + topic\_id\\score, evidence JSON};
  \node[entity, right=of papertopic] (authortopic) {\textbf{Canonical Author / Alias / Topic}\\identity and merge state\\publication link, paper count\\score, evidence JSON};

  \draw[arrow] (paper) -- node[above,font=\scriptsize]{1:N} (chunk);
  \draw[dashedarrow] (chunk) -- node[above,font=\scriptsize]{IDs in JSON} (answer);
  \draw[arrow] (paper) -- node[left,font=\scriptsize]{1:N} (artifact);
  \draw[dashedarrow] (chunk) -- node[right,font=\scriptsize]{retrieved evidence} (recommendation);
  \draw[dashedarrow] (answer) -- node[right,font=\scriptsize]{polymorphic item} (review);
  \draw[arrow] (paper) |- (papertopic);
  \draw[arrow] (topic) -- (papertopic);
  % Route below the middle association node.  An explicit waypoint avoids the
  % ambiguous border-intersection route that PGF reports when ``-|'' crosses
  % the PaperTopic node before entering AuthorTopic.
  \draw[arrow] (topic.south) |- ($(papertopic.south)+(0,-5mm)$) -| (authortopic.south);
  \draw[dashedarrow] (artifact) -- node[left,font=\scriptsize]{may supply inferred topic text} (topic);
\end{tikzpicture}
}
\caption[Principal persistence and provenance entities.]{Principal persistence and provenance entities. Solid arrows denote ownership/foreign-key relations; dashed arrows denote evidence or polymorphic references.}
\label{fig:entity-model}
\end{figure}

\section{Authoritative index design}

Each vector manifest binds the corpus and ordered chunks/hashes to provider/model/revision, dimension/normalisation, configuration/code/index hashes, role, and completeness. A rebuild validates an immutable generation, commits matching database status, then atomically moves the current pointer; failure leaves the prior pointer active. Demo indexes cannot target the canonical location, and validators reject missing, partial, stale, or mismatched state.

Feature hashing is a 256-dimensional signed lexical projection, not learned semantics. Dense retrieval pins \texttt{all-MiniLM-L6-v2} revision \texttt{826711e...a1def9}, its artifact hash, 384 dimensions, normalization, CPU/window pooling, and license. Production Ollama use requires an approved model/digest. An explicit local-demo policy may permit installed models while still checking the observed digest before and after generation; absent response attestation is recorded as a mutable-tag limitation.

\section{Security and privacy boundaries}

Public reads/generation are separate from technical evaluation and protected mutations. Production metadata requires human approval, published state and cleared rights; searchable text also requires approved extraction/current public generation. No audited record met the full predicate, so production evidence routes fail closed. Local administrators authenticate with Argon2id passwords and HttpOnly, SameSite sessions; a readable CSRF cookie must match the mutation header. Idle/absolute expiry, lockout, forced temporary-password replacement, role restrictions, production HTTPS/host/CORS checks, size/rate limits, redaction, SSRF/redirect controls, and PDF limits address principal threats. Environment-digested bearer service actors remain available for bounded automation.

Student interests, skills, time, constraints, topic preferences, and Idea Generator conversation turns remain client-side/in-request and are not persisted by default. Public generated-output serializers omit private reviewer notes and draft corrections. This minimisation reduces privacy risk but does not replace an institutional privacy review for a deployed service.
